\documentclass[10pt,a4paper]{amsart}
\usepackage[margin=19mm]{geometry}
\usepackage{amsmath,amssymb,mathtools}
\usepackage[scr=rsfs,cal=euler]{mathalfa}
\usepackage{xfrac}
\usepackage[unicode,hidelinks,bookmarks=false]{hyperref}
\newcommand{\ie}{i.e.\ }
\newcommand{\cf}{cf.\ }
\newcommand{\resp}{respectively\ }
\newcommand{\Subsection}{\subsection}
\providecommand{\nobreakdash}{\nobreak\hbox{-}}
\setlength{\emergencystretch}{2em}
\multlinegap0pt
\usepackage{bm}
\usepackage{bbm}
\usepackage{dsfont}
\usepackage{xspace}
\usepackage{cancel}
\usepackage{mathtools}
\usepackage{amsthm}
\makeatletter
\@ifundefined{theorem}{\newtheorem{theorem}{Theorem}}{}
\@ifundefined{lemma}{\newtheorem{lemma}[theorem]{Lemma}}{}
\@ifundefined{proposition}{\newtheorem{proposition}[theorem]{Proposition}}{}
\makeatother
\DeclareMathOperator{\diam}{diam}
\DeclareMathOperator{\lex}{lex}
\DeclareMathOperator{\mingens}{Mingens}
\title[Trees whose path ideals have linear quotients]{Trees whose path ideals have linear quotients}
\author{Trung Chau, Kanoy Kumar Das, Animikha Dutta Dhar, Pranath S Karanth, Aniruda Suswaram}
\date{Source: arXiv:2506.06209v1 (CC BY 4.0)}
\begin{document}
\maketitle

\noindent\emph{Source and licence.} Trees whose path ideals have linear quotients. Version arXiv:2506.06209v1 (CC BY 4.0), distributed under the Creative Commons Attribution 4.0 International licence, as recorded in the arXiv licence metadata for this exact version and shown on the abstract page https://arxiv.org/abs/2506.06209v1. Full rights notice: licenses/arxiv-2506.06209-CC-BY-4.0.txt.\par\smallskip

\noindent\textit{Editorial scope.} Counted unit: the Proposition whose source label is \texttt{prop:LQ-2n-3} in the article (source0.tex lines 882--884, source label \texttt{prop:LQ-2n-3}), delivered together with the article's own complete original proof of it (source0.tex lines 885--928, one \texttt{proof} environment of 44 lines). No mathematics is written, repaired or re-derived: statement and proof are the article's own text line for line. Required context retained uncounted: lem:gens-Jn-caterpillar (lines 873--875). Every definition, display and label of those lines is retained unchanged. Omitted from this package and not redistributed: 1--872, 876--881, 929--1061. Nothing inside a retained excerpt is omitted or altered: every retained body is delivered exactly as the article's lines are written, so no omission receipt of any kind accompanies this package. Citations: the retained excerpts carry no citation command at all, so no bibliography entry of the article is transcribed and no reference is redistributed. Reference adaptation: none was needed and none was made; every reference inside the delivered unit resolves inside it, so no unresolved reference or citation is produced. Printed slips kept literal: no printed word, symbol, hypothesis or number of the counted statement or of its proof was altered. Layout and class: the article's own class is \texttt{amsart} with packages including geometry, xcolor, amsfonts, amsthm, amssymb, verbatim, amscd, amsmath, blindtext, multirow, multicol, graphicx, enumitem, yhmath; the wrapper is the lane's pinned \texttt{amsart} editorial wrapper at 10pt on A4, which restates the theorem-like environments the retained text needs and adds the article's own macro definitions that the retained text needs (\texttt{\textbackslash diam}, \texttt{\textbackslash lex}, \texttt{\textbackslash mingens}), with extra wrapper packages (bm, bbm, dsfont, xspace, cancel, mathtools). The delivered numbering is the unit's own. Assets: no retained span contains a figure environment, a graphics-inclusion command, a PDF-page inclusion, a table or a TikZ picture; no image file is copied into this package, the package declares no asset dependency and no image file is redistributed with it. Licence: version arXiv:2506.06209v1 of this article is distributed under the Creative Commons Attribution 4.0 International licence, as recorded in the arXiv licence metadata for this exact version and shown on the abstract page https://arxiv.org/abs/2506.06209v1. A full rights notice is delivered at licenses/arxiv-2506.06209-CC-BY-4.0.txt. Characters: every retained excerpt is pure ASCII, so the delivered engine, XeLaTeX with its default Latin Modern text font, reports no missing character.
	\begin{lemma}\label{lem:gens-Jn-caterpillar}
		Let $G$ be a caterpillar tree and $n\geq 2$ a positive integer. Then the minimal monomial generators of $J_n(G)$ are of the form $yx_ix_{i+1}\cdots x_{i+n-3}z$ where $i\in [1,\ d-n+2]$, $y\in LN_G(x_i)\cup \{x_{i-1}\}$, and $z\in  LN_G(x_{i+n-3})\cup \{x_{i+n-2}\}$.
	\end{lemma}

	\begin{proposition}\label{prop:LQ-2n-3}
		Let $G$ be a caterpillar tree and $n\geq 4$ a positive integer. Assume that $d=\diam(G)\leq 2n-3$. Then $J_n(G)$ has linear quotients with respect to $(>_{\lex})$.
	\end{proposition}

	\begin{proof}
		Let $m\in \mingens(J_n(G))$. By Lemma~\ref{lem:gens-Jn-caterpillar}, we can set \[m=yx_ix_{i+1}\cdots x_{i+n-3}z\]
		where $i\in [1,\ d-n+2]$, $y\in LN_G(x_i)\cup \{x_{i-1}\}$, and $z\in  LN_G(x_{i+n-3})\cup \{x_{i+n-2}\}$. It now suffices to show that $J_n(G)_{>_{lex}m} \colon m$ is generated by variables, so long as $m$ is not the biggest monomial in $\mingens(J_n(G))$ with respect to $(>_{\lex})$. We have two cases.
		
		\textbf{Case 1:} Assume that $m=x_{i-1}x_ix_{i+1}\cdots x_{i+n-3}x_{i+n-3,v}$ where $i\in [2, d-n+2]$ and $v\in [1, l_{i+n-3}+1]$. For any $y\in LN_G(x_{i-1}) \cup \{x_{i-2}\} \cup \{x_{i+n-3,k}\colon k\in [1,v)\}$, we have $y\succ x_{i+n-3,v}$. In particular, this means that  $m\frac{y}{x_{i+n-3,v}} >_{\lex}  m$, and it is in $\mingens(J_n(G))$ as this monomial corresponds to a path of length $n-1$. Therefore, we have 
		\[
		(LN_G(x_{i-1}) \cup \{x_{i-2}\} \cup \{x_{i+n-3,k}\colon k\in [1,v)\}) \subseteq (J_n(G)_{>_{lex}m}\colon m).
		\]
		We want to show that the converse also holds. Indeed, consider $m'=yx_jx_{j+1}\cdots x_{j+n-3}z>_{\lex}  m$ where $j\in [1,\ d-n+2]$, $y\in LN_G(x_j)\cup \{x_{j-1}\}$, and $z\in  LN_G(x_{j+n-3})\cup \{x_{j+n-2}\})$. It now suffices to show that $(m':m)\subseteq (LN_G(x_{i-1}) \cup \{x_{i-2}\} \cup \{x_{i-n+3,k}\colon k\in [1,v)\})$.
		
		Due to the fact that $m'>_{\lex}  m$ and the structure of $m$, we must have $j<i$, or $j=i, y=x_{j-1}$ and $z\in \{x_{i-n+3,k}\colon k\in [1,v)\}$. Clearly if $j=i$, then $(m'\colon m)\subseteq (z)\subseteq (\{x_{i-n+3,k}\colon k\in [1,v)\})$, and thus the result follows. Note that $z$ here makes sense if and only if $v\geq 2$, and we can assume this as when $v=1$, $m$ is the biggest monomial in $\mingens(J_n(G))$ with respect to $(>_{\lex})$.
		
		We can now assume that $j<i$. If $i=2$ then $j=1$, and we have $m=x_1x_2\cdots x_{n-1}x_{n-1,v}$ (implicitly, this implies that $n\leq d-1$, but we will not need this), and $m'=yx_1x_2\cdots x_{n-2} z$ where $y\in LN_G(x_1)$. Hence $(m'\colon m) \subset (y)\in (LN_G(x_1))=(LN_G(x_{i-1}))$, as desired. Now we can assume that $i\geq 3$, i.e., $x_{i-2}$ is a vertex in $\{x_1,\dots, x_{d-1}\}$. We then have three subcases:
		\begin{itemize}
			\item Assume that $i-2\in [j,j+n-3]$. Then $x_{i-2}\mid m'$, and thus $(m'\colon m) \subseteq (x_{i-2})$, as~desired.
			\item Assume that $i-2<j$. Recall that we have $j<i$. Thus $j=i-1$. Therefore $(m'\colon m)\subseteq (y)\subseteq (LN_G(x_{j})\cup \{x_{j-1}\}) = (LN_G(x_{i-1})\cup \{x_{i-2}\})$, as desired.
			\item Assume that $i-2\geq j+n-3$. Then \[
			j<i-n+1 \leq (d-n+2)-n+1 = d-2n+3 \leq (2n-3)-2n+3 =0,
			\]
			a contradiction.
		\end{itemize}
		
		\textbf{Case 2:} Assume that $m=x_{i,u}x_ix_{i+1}\cdots x_{i+n-3}x_{i+n-3,v}$ where $i\in [1, d-n+2]$, $u\in [1,l_i]$, and $v\in [1, l_{i+n-3}+1]$.    
		For any $y\in \{ x_{i,k} \colon k\in [1,u) \} \cup \{x_{i-1}\} \cup \{x_{i+n-3,k}\colon k\in [1,v)\}$, we have $y\succ x_{i+n-3,v}$. In particular, this means that  $m\frac{y}{x_{i+n-3,v}} >_{\lex} m$, and it is in $\mingens(J_n(G))$ as this monomial corresponds to a path of length $n-1$. Therefore, we have 
		\[
		(\{ x_{i,k} \colon k\in [1,u) \} \cup \{x_{i-1}\} \cup \{x_{i+n-3,k}\colon k\in [1,v)\}) \subseteq (J_n(G)_{>_{lex}m} \colon m).
		\]
		We want to show that the converse also holds. Indeed, consider $m'=yx_jx_{j+1}\cdots x_{j+n-3}z>_{\lex}  m$ where $j\in [1,\ d-n+2]$, $y\in LN_G(x_j)\cup \{x_{j-1}\}$, and $z\in  LN_G(x_{j+n-3})\cup \{x_{j+n-2}\}$. It now suffices to show that $(m':m)\subseteq (\{ x_{i,k} \colon k\in [1,u) \} \cup \{x_{i-1}\} \cup \{x_{i+n-3,k}\colon k\in [1,v)\})$.
		
		Due to the fact that $m'>_{\lex}  m$ and the structure of $m$, one of the following holds:
		\begin{enumerate}
			\item[(a)] $j<i$;
			\item[(b)] $j=i$ and $y\in \{ x_{i,k}\colon k\in [1,u)\}$;
			\item[(c)] $j=i, y=x_{i,u}$, and $z\in \{x_{i-n+3,k}\colon k\in [1,v)\}$. 
		\end{enumerate}
		If (b) holds, then $(m'\colon m)\subseteq (y)\subseteq (\{ x_{i,k}\colon k\in [1,u)\})$, and thus the result follows. Similar arguments apply when (c) holds. We can now assume that (a) holds, i.e., $j\leq i-1$. We have 
		\begin{align*}
			(i-1)-(j+n-3) = i-n+2-j &\leq (d-n+2) -n+2-(1) \\
			&= d-2n+3\\
			&\leq (2n-3)-2n+3\\
			&=0.        
		\end{align*}
		Thus, $i-1\in [j,j+n-3]$. Therefore, $(m'\colon m)\subseteq (x_{i-1}),$ as desired. 
	\end{proof}

\end{document}
